
\documentclass{article}
\usepackage{colortbl}
\usepackage{makecell}
\usepackage{multirow}
\usepackage{supertabular}

\begin{document}

\newcounter{utterance}

\centering \large Interaction Transcript for game `cladder', experiment `full\_v1.5\_default', episode 1876 with llama{-}3.1{-}8b{-}instruct.
\vspace{24pt}

{ \footnotesize  \setcounter{utterance}{1}
\setlength{\tabcolsep}{0pt}
\begin{supertabular}{c@{$\;$}|p{.15\linewidth}@{}p{.15\linewidth}p{.15\linewidth}p{.15\linewidth}p{.15\linewidth}p{.15\linewidth}}
    \# & \multicolumn{2}{c}{Player} && \multicolumn{2}{c}{Game Master} \\
    \hline

    \theutterance \stepcounter{utterance}  
    & & & \multicolumn{4}{p{0.6\linewidth}}{
        \cellcolor[rgb]{0.9,0.9,0.9}{
            \makecell[{{p{\linewidth}}}]{
                \texttt{\tiny{[P1$\langle$GM]}}
                \texttt{You are an expert in causal inference. The following question is not a typical commonsense query, but rather a meticulously designed question created by a professor specializing in causal inference, intended to assess the students' mastery of the course content.} \\
\texttt{Imagine a self{-}contained, hypothetical world with only the following conditions, and without any unmentioned factors or causal relationships: Medication has a direct effect on blood pressure and heart condition. Blood pressure has a direct effect on heart condition. For people not taking any medication, the probability of healthy heart is 64\%. For people taking medication, the probability of healthy heart is 21\%. For people taking medication, would it be more likely to see healthy heart if the person had taken no medication?} \\
\texttt{Start your answer with 'Yes' or 'No', followed by additional reasoning or evidence to support your explanation.} \\
            }
        }
    }
     \\ \\

    \theutterance \stepcounter{utterance}  
    & \multicolumn{4}{p{0.6\linewidth}}{
        \cellcolor[rgb]{0.561,0.737,0.561}{
            \makecell[{{p{\linewidth}}}]{
                \texttt{\tiny{[P1$\rangle$GM]}}
                \texttt{No.} \\
\\ 
\texttt{To answer this question, we need to consider the causal relationships and the conditional probabilities given. We are given that:} \\
\\ 
\texttt{1. Medication has a direct effect on blood pressure and heart condition.} \\
\texttt{2. Blood pressure has a direct effect on heart condition.} \\
\texttt{3. For people not taking any medication, the probability of healthy heart is 64\%.} \\
\texttt{4. For people taking medication, the probability of healthy heart is 21\%.} \\
\\ 
\texttt{We want to know whether it is more likely to see a healthy heart if the person had taken no medication, given that they are taking medication. This is a counterfactual question, which can be addressed using the potential outcomes framework.} \\
\\ 
\texttt{Let's denote the potential outcomes as follows:} \\
\\ 
\texttt{{-} H: Healthy heart} \\
\texttt{{-} UH: Unhealthy heart} \\
\texttt{{-} M: Taking medication} \\
\texttt{{-} NM: Not taking medication} \\
\\ 
\texttt{We are interested in the probability of H given M, which we can denote as P(H|M). We can also consider the probability of H given NM, which we can denote as P(H|NM).} \\
\\ 
\texttt{Using the law of total probability, we can express P(H|M) as:} \\
\\ 
\texttt{P(H|M) = P(H|M,NM) \textbackslash{}* P(NM|M) + P(H|M,M) \textbackslash{}* P(M|M)} \\
\\ 
\texttt{Since M and NM are mutually exclusive, P(NM|M) = 0. Therefore, the equation simplifies to:} \\
\\ 
\texttt{P(H|M) = P(H|M,M) \textbackslash{}* P(M|M)} \\
\\ 
\texttt{We are given that P(H|M) = 21\%. We can also calculate P(M|M) = 1, since the person is taking medication.} \\
\\ 
\texttt{However, to answer the question, we need to consider the probability of H given that the person had taken no medication, which is P(H|NM,M). This is a counterfactual quantity, which cannot be directly observed. We can use the potential outcomes framework to express this quantity as:} \\
\\ 
\texttt{P(H|NM,M) = P(H|M,NM) / P(NM|M)} \\
\\ 
\texttt{Since P(NM|M) = 0, this expression is undefined. However, we can use the law of total probability to express P(H|M,NM) as:} \\
\\ 
\texttt{P(H|M,NM) = P(H|NM) \textbackslash{}* P(NM|NM)} \\
\\ 
\texttt{We are given that P(H|NM) = 64\%. We can also calculate P(NM|NM) = 1, since the person is not taking medication.} \\
\\ 
\texttt{Therefore, P(H|M,NM) = 64\%. However, this is not the probability we are interested in. We want to know whether it is more likely to see a healthy heart if the person had taken no medication, given that they are taking medication. This is equivalent to comparing P(H|M,NM) and P(H|M).} \\
\\ 
\texttt{Since P(H|M,NM) = 64\% and P(H|M) = 21\%, it is clear that P(H|M,NM) > P(H|M).} \\
            }
        }
    }
    & & \\ \\

    \theutterance \stepcounter{utterance}  
    & & & \multicolumn{2}{p{0.3\linewidth}}{
        \cellcolor[rgb]{0.9,0.9,0.9}{
            \makecell[{{p{\linewidth}}}]{
                \texttt{\tiny{[GM$|$GM]}}
                \texttt{No} \\
            }
        }
    }
    & & \\ \\

    \theutterance \stepcounter{utterance}  
    & & & \multicolumn{2}{p{0.3\linewidth}}{
        \cellcolor[rgb]{0.9,0.9,0.9}{
            \makecell[{{p{\linewidth}}}]{
                \texttt{\tiny{[GM$|$GM]}}
                \texttt{yes} \\
            }
        }
    }
    & & \\ \\

    \theutterance \stepcounter{utterance}  
    & & & \multicolumn{2}{p{0.3\linewidth}}{
        \cellcolor[rgb]{0.9,0.9,0.9}{
            \makecell[{{p{\linewidth}}}]{
                \texttt{\tiny{[GM$|$GM]}}
                \texttt{game\_result = LOSE} \\
            }
        }
    }
    & & \\ \\

\end{supertabular}
}

\end{document}
